\begin{textsample}{POS Dim 2 – human – Score 6.00 – mg/mg\_ef\_linguagens\_ed\_fis\_1\_p4.txt}  \label{ex:f2_pos_014}
A materialidade de nossa existência está sustentada pela nossa corporeidade : do momento da concepção até nosso último minuto de vida estamos vinculados aos movimentos cíclicos, acíclicos, voluntários, involuntários, conscientes, expressivos, etc. de partes do nosso corpo ou dele como um todo.
Produzimos valores, \textbf{saberes}, conhecimentos, princípios, normas e conceitos a \textbf{partir} das expressões e relações de todas as nossas dimensões constitutivas.
Nossa \textbf{sociedade} ainda valoriza muito a dimensão mental/cognitiva como processo de aquisição e produção de conhecimento.
Mas diversas teorias têm oferecido legitimidade e importância para a participação de outras dimensões – motora/corporal, afetivo/emocional, social/cultural e espiritual/existenciais- neste processo.
Ou ainda mais, estas dimensões têm sido validadas como imprescindíveis para a formação \textbf{humana} e cidadã atualmente almejada, também expressa nas 10 Competências Gerais da Base Nacional Comum Curricular.
Respaldamos neste sentido as aulas de Educação Física na escola como tempo e espaço privilegiado que pode garantir o reconhecimento e a valorização da vivência destas dimensões e seu potencial formativo.

% matched lemmas: humano, partir, saber, sociedade
\end{textsample}
